% 01-problem.tex: the problem Polaris is trying to solve.
\section{The problem}
\label{sec:problem}

A person in the United States carries six to eight credentials that do not talk to each other:
a driver's licence, a passport, a Social Security card, a Real ID, a voter registration, an
insurance card. Each is a separate artifact, signed by a separate authority to a separate standard,
with no shared revocation path and no shared audit trail. Version 1 of this report set out to model
their consolidation into one active credential record per person, verified through context-scoped
events. That is the easy half of the problem, and it is still the problem's shape.

The hard half is what happens when the system works. An identity system that succeeds becomes the
doorway to everything, and a doorway is a place where power concentrates. Three failure modes
follow, and every layer of Polaris exists to answer one of them.

\begin{description}
  \item[The system lies about the past.] An operator, or an insider with database access,
    revokes a token and later erases the record; issues a credential under procedures the public
    was never told of; or shows one history to an auditor and another to a court. A national
    identity system cannot be allowed to retroactively rewrite what it did.
  \item[The system watches.] Every verification is a fact about where a person was and what they
    were doing. A log of those facts is a surveillance backbone whether or not anyone intended one,
    and it stays one after the operator's intentions change.
  \item[The system coerces.] A holder can be compelled: a border officer, an employer, or a
    thief can force a credential to be presented, and a doorway that every service requires can
    force a population to enrol. An identity system is also the most efficient instrument of
    exclusion ever built, if the list of who is not in it can be produced on demand.
\end{description}

The conventional answer to all three is policy: rules about retention, access, audit and consent,
enforced by the people who run the system. Polaris's constitution (\code{MISSION.md}) takes the
position that policy is worthless as the only thing holding a promise up, because the people who
run the system are exactly the ones the promise is meant to bind. The promise has to be enforced
where it cannot be talked around: in the schema, in triggers, in unique indexes, in check
constraints, in signatures that anyone can verify without asking the system, and in logs that other
people watch. The application is not the security boundary. The database is, and the layers around
the database exist to make even the database's claims checkable by strangers.

\subsection{Why post-quantum from the first day}

Version 1 argued, through Mosca's inequality, that identity data has a secrecy horizon measured in
generations, that migration of a national system takes a decade, and that a credential system
deployed on classical cryptography in 2026 is therefore a deferred compromise rather than a working
system. That argument is unchanged and is summarised in Appendix~\ref{app:from-v1}. At \code{v9.345}
the default signing algorithm is ML-DSA-65 (FIPS 204), ML-DSA-87 is accepted beside it, the zero-knowledge
proof system is hash-based with no trusted setup, and the public TLS edge negotiates a hybrid
post-quantum key exchange with modern clients. Which surfaces remain classical is stated plainly in
Section~\ref{sec:operations} and in the repository's posture ledger.

\subsection{The progression this paper follows}

Each of the following sentences names a layer, and each is true of the layer before it.

\begin{itemize}
  \item A credential's authenticity is not its current authorization. A signature stays genuine
    forever; the authority under it can be withdrawn this afternoon.
  \item Current authorization is not privacy. A verifier can learn that a credential is valid and
    still learn far more than it needed to.
  \item Privacy is not freedom from coercion. A system can know almost nothing about a person and
    still be the doorway every service demands.
  \item Federation is not institutional independence. Two instances of the same code run by the
    same author prove that the protocol works, not that independent institutions run it.
  \item A transparency log is not protection from a split view. One observer can be shown a private
    fork of history and never know.
  \item A green test is not proof that the test asked the right question.
\end{itemize}

The last sentence is the one the project has learned most expensively, and Section~\ref{sec:senses}
is about it.
